% Source: 16_The_Age_of_the_Christian_Roman_Empire_312_590.txt
\bookpart[The Age of the Christian Roman Empire]{THE AGE OF THE CHRISTIAN ROMAN EMPIRE}{312–590}

\begin{partsummary}
The Emperor Constantine is one of the major figures of Christian history. After his conversion
Christianity moved swiftly from the seclusion of the catacombs to the prestige of palaces. The
movement started the fourth century as a persecuted minority; it ended the century as the
established religion of the empire. Thus the Christian church was joined to the power of the state
and assumed a moral responsibility for the whole society. Initially under the instruction of
Constantine, the church refined its doctrine and developed its structure. Some, such as the
historian Eusebius, saw Constantine’s embrace of Christianity as its victory over the empire.
Others, such as the monks, believed the culture was capturing Christianity. The story that follows
is the “Christianizing” of the great Hellenistic mindset and world. When the empire gave way to
barbarian invaders (known as Europeans today), the monks ironically won over the conquerors; the
monks displayed the dignity of a well-ordered life in community with deep roots in the Christian
faith.
\end{partsummary}

\timelinetitle{The Age of the Christian Roman Empire}
